\section{What Mitigates Failures, and at What Cost}\label{sec:mitigation}
The previous sections asked where failures are decided. This one collects what the same data say
about mitigating them cheaply (Fig.~\ref{fig:mitigation}, Table~\ref{tab:mitigation}).

\textbf{Maths: restarting recovers most failures.} For the failed ProcessBench solutions of the four
generators, plain restarting (independent attempts from the question, FLOP-matched, every token
counted) solves \recMinOne--\recMaxOne\% of them with one attempt, \recMinFour--\recMaxFour\% within
four and \recMinEight--\recMaxEight\% within eight. No surgical policy in Table~\ref{tab:policies}
does better at any of these budgets, except PRM backtracking for Qwen2.5-Math-7B at eight attempts
(\QwenMathPolPrmDropEightDiff, exploratory).

\textbf{RAG: detect, restart, abstain.} The detector $D$ is a 150M-parameter token classifier, so
checking a draft costs little; only flagged drafts are resampled, and what is still flagged after
the budget is withheld (QA) or deleted (other tasks). This loop (DRA, ``restart until $D$ passes'' in
Section~\ref{sec:grounded}) is the cheapest mitigation that resamples: \tokDra\ tokens per item against
\tokBase\ for the base model (\costDra$\times$), while claim repair costs \costRep$\times$, the
choke-point-aware controller \costCare$\times$ and Chain-of-Verification \costCove$\times$. Deleting
flagged sentences costs almost nothing (\costDel$\times$) but changes the content of the answer;
abstaining on flagged QA gives the lowest hallucination under every grader at a coverage of
\covAbst\%.

\begin{table}[t]\centering\small
\caption{Mitigation on RAGTruth (450 sources, Llama-2-7B-chat). Tokens per item, every call counted;
hallucination rate (\%) under the independent judge $E_1$ and the detector-family grader $E_2$.
CoVe: raw outputs on 192 items (the echoed verification block inflates both graders,
Section~\ref{sec:grounded}).}
\label{tab:mitigation}
\setlength{\tabcolsep}{3pt}
\resizebox{\columnwidth}{!}{%
\begin{tabular}{lccccc}\toprule
Mitigation & Tokens & Cost & $E_1$ & $E_2$ & Coverage\\\midrule
None (base model) & \tokBase & 1.00$\times$ & \hallEOneBase & \hallETwoBase & \covBase\\
Delete flagged sentences & \tokDel & \costDel$\times$ & \hallEOneDel & \hallETwoDel & \covDel\\
\textbf{Detect, restart, abstain (DRA)} & \tokDra & \costDra$\times$ & \hallEOneDra & \hallETwoDra & \covDra\\
Detect, backtrack, abstain & \tokBack & \costBack$\times$ & \hallEOneBack & \hallETwoBack & \covBack\\
Claim repair, then delete & \tokRep & \costRep$\times$ & \hallEOneRep & \hallETwoRep & \covRep\\
CARE (choke-point-aware) & \tokCare & \costCare$\times$ & \hallEOneCare & \hallETwoCare & \covCare\\
Abstain on flagged QA & \tokAbst & \costAbst$\times$ & \hallEOneAbst & \hallETwoAbst & \covAbst\\
Chain-of-Verification & \tokCove & \costCove$\times$ & \hallEOneCove & \hallETwoCove & \covCove\\\bottomrule
\end{tabular}}
\end{table}

\begin{figure*}[t]\centering
\includegraphics[width=0.95\textwidth]{fig27_mitigation.pdf}
\caption{What mitigation buys. Left: share of failed maths solutions solved by plain restarting, by
budget (95\% CIs). Right: hallucination after each RAG mitigation under the independent judge
($E_1$) and the detector-family grader ($E_2$), with the human base rate and the human-scale
prediction of Proposition~3 (Appendix~\ref{app:more}) for DRA; cost relative to no mitigation in brackets.}
\label{fig:mitigation}
\end{figure*}

\textbf{How much does DRA reduce hallucination for people?} The two graders disagree by an order of
magnitude, and no human rating of system outputs exists yet. Proposition~3 (Appendix~\ref{app:more}) gives a
principled estimate. If the system returns the first independent draw that $D$ passes, the share
of returned answers that are hallucinated is $P(\text{hallucinated}\mid D\text{ passes})$, whatever
the retry budget. RAGTruth's human labels on the drafts measure exactly this quantity: $D$ misses
\lawBeta\% of hallucinated drafts and flags \lawAlpha\% of faithful ones, so the predicted rate among
DRA's answers is \lawPi\% [\lawPiLo, \lawPiHi], against \lawH\% [\lawHLo, \lawHHi] for the base model:
a reduction of \lawRed\ points [\lawRedLo, \lawRedHi]. Two assumptions carry this estimate (draws
independent given the input, and the detector's error rates on resampled answers equal those on
drafts), and with per-item variation in hallucination rates it is optimistic, so we read it as an
upper bound on the reduction and report it as exploratory. An attempt to confirm it by inverting the
independent judge's own error rates was consistent (\lawInv\%) but too imprecise to be informative
([\lawInvLo, \lawInvHi]); the human rating kit is released for that purpose.

\textbf{What the law says about cost.} The retry budget only buys coverage: abstentions fall
geometrically with more attempts, but the hallucination rate of delivered answers stays at the
detector's floor. The lever is therefore the detector, not the budget and not the resampling point:
a better detector lowers the floor, and restarting is the cheapest way to cash it in.
